%% ════════════════════════════════════════════════════════════════════════════
\section{Full REST API Endpoint Reference}
\label{app:api}
%% ════════════════════════════════════════════════════════════════════════════

Table~\ref{tab:app_api} lists all REST endpoints with method, path,
authentication requirement, and role restriction.

\begin{table}[H]
\centering
\caption{Complete REST API Endpoint Reference}
\label{tab:app_api}
\begin{tabular}{@{}p{0.8cm}p{3.6cm}p{1.2cm}p{2.0cm}@{}}
\toprule
\textbf{M.} & \textbf{Path} & \textbf{Auth} & \textbf{Role} \\
\midrule
POST & \texttt{/api/auth/login}               & No  & Any \\
GET  & \texttt{/api/sensors/latest}           & JWT & Any \\
GET  & \texttt{/api/sensors/history}          & JWT & Any \\
GET  & \texttt{/api/sessions/\{id\}/sensors/latest} & JWT & Any \\
GET  & \texttt{/api/sessions/\{id\}/sensors/summary} & JWT & Any \\
POST & \texttt{/api/sessions/start}           & JWT & Prof \\
POST & \texttt{/api/sessions/\{id\}/end}      & JWT & Prof \\
GET  & \texttt{/api/sessions}                 & JWT & Prof \\
GET  & \texttt{/api/sessions/\{id\}}          & JWT & Prof \\
POST & \texttt{/api/control/ac}               & JWT & Prof \\
POST & \texttt{/api/control/lighting}         & JWT & Prof \\
GET  & \texttt{/api/control/status/\{room\}}  & JWT & Prof \\
GET  & \texttt{/api/alerts}                   & JWT & Prof \\
PATCH& \texttt{/api/alerts/\{id\}/acknowledge}& JWT & Prof \\
GET  & \texttt{/api/alerts/unread-count/\{room\}} & JWT & Prof \\
GET  & \texttt{/api/courses}                  & JWT & Prof \\
POST & \texttt{/api/courses}                  & JWT & Prof \\
POST & \texttt{/api/courses/\{id\}/enroll}    & JWT & Prof \\
GET  & \texttt{/api/sessions/\{id\}/attendance} & JWT & Prof \\
PATCH& \texttt{/api/attendance/\{id\}}        & JWT & Prof \\
POST & \texttt{/api/sessions/\{id\}/mark-absent} & JWT & Prof \\
GET  & \texttt{/api/students}                 & JWT & Prof \\
POST & \texttt{/api/students}                 & JWT & Prof \\
POST & \texttt{/api/students/\{id\}/enroll-face} & JWT & Prof \\
GET  & \texttt{/api/at-risk}                  & JWT & Prof \\
GET  & \texttt{/api/at-risk/\{student\_id\}}  & JWT & Prof \\
POST & \texttt{/api/at-risk/recompute}        & JWT & Admin \\
GET  & \texttt{/api/forecasting}              & JWT & Prof \\
GET  & \texttt{/api/forecasting/\{course\_id\}} & JWT & Prof \\
POST & \texttt{/api/forecasting/recompute}    & JWT & Admin \\
GET  & \texttt{/api/webcam/recognition-status} & JWT & Prof \\
POST & \texttt{/api/webcam/start-recognition} & JWT & Prof \\
POST & \texttt{/api/webcam/stop-recognition}  & JWT & Prof \\
GET  & \texttt{/api/webcam/encodings}         & JWT & Prof \\
POST & \texttt{/api/webcam/attendance}        & JWT & Prof \\
WS   & \texttt{/ws/classroom/\{room\_id\}}    & No  & — \\
\bottomrule
\end{tabular}
\end{table}

%% ════════════════════════════════════════════════════════════════════════════
\section{MQTT Topic Schema}
\label{app:mqtt}
%% ════════════════════════════════════════════════════════════════════════════

All topics follow the pattern \texttt{classroom/\{room\_id\}/\{suffix\}}.
The default \texttt{room\_id} is \texttt{room1}, configurable via the
\texttt{ROOM\_ID} environment variable and the ESP32 \texttt{config.h}.

\begin{table}[H]
\centering
\caption{MQTT Full Topic Schema with JSON Payload Schemas}
\label{tab:app_mqtt}
\begin{tabular}{@{}p{2.5cm}p{1.0cm}p{4.0cm}@{}}
\toprule
\textbf{Suffix} & \textbf{QoS} & \textbf{JSON Schema} \\
\midrule
\texttt{sensors/temperature} & 0 & \texttt{\{"value": float, "unit": "C", "ts": int\}} \\
\texttt{sensors/humidity}    & 0 & \texttt{\{"value": float, "unit": "\%", "ts": int\}} \\
\texttt{sensors/air\_quality}& 0 & \texttt{\{"value": int, "unit": "ppm", "ts": int\}} \\
\texttt{sensors/sound}       & 0 & \texttt{\{"value": 0|1, "unit": "bool", "ts": int\}} \\
\texttt{status}              & 0 & \texttt{\{"online": true, "ts": int\}} \\
\texttt{relay/ac}            & 0 & \texttt{\{"action": "on"|"off"|"auto"\}} \\
\texttt{relay/lighting}      & 0 & \texttt{\{"action": "on"|"off"|"auto"\}} \\
\texttt{alerts}              & 0 & \texttt{\{"type": str, "value": float\}} \\
\bottomrule
\end{tabular}
\end{table}

%% ════════════════════════════════════════════════════════════════════════════
\section{Database Schema}
\label{app:db}
%% ════════════════════════════════════════════════════════════════════════════

The following definitions correspond to the SQLAlchemy ORM models in
\texttt{backend/app/models/db\_models.py}. All primary keys are UUIDs
generated server-side.

\begin{verbatim}
professors(id PK, name, email UNIQUE, hashed_password,
           role ENUM(professor|admin), created_at)

students(id PK, name, student_id UNIQUE, created_at)

courses(id PK, code UNIQUE, name,
        professor_id FK→professors NULLABLE,
        professor_name)

sessions(id PK, course_id FK→courses, room_id,
         started_at, ended_at NULLABLE,
         status ENUM(active|ended|upcoming))

attendance_records(id PK, session_id FK, student_id FK,
  status ENUM(present|absent|late|excused),
  detected_at, adjusted_by NULLABLE,
  adjusted_at NULLABLE, moodle_synced BOOL DEFAULT false)

face_encodings(id PK, student_id FK→students,
               encoding BYTEA, created_at)

sensor_readings(id PK, room_id, sensor_type ENUM,
                value FLOAT, unit, recorded_at)

alerts(id PK, room_id, type ENUM, value FLOAT NULLABLE,
       message TEXT, acknowledged BOOL DEFAULT false,
       created_at)

at_risk_explanations(id PK, student_id FK UNIQUE,
  overall_attendance_rate FLOAT,
  summary_explanation TEXT,
  per_course_data JSONB,
  generated_at TIMESTAMP,
  ollama_reachable BOOL)

attendance_forecasts(id PK, course_id FK UNIQUE,
  trend_data JSONB,
  expected_next_rate FLOAT NULLABLE,
  trend_classification VARCHAR(30) NULLABLE,
  confidence_level VARCHAR(10) NULLABLE,
  interpretation TEXT NULLABLE,
  suggested_action VARCHAR(30) NULLABLE,
  ollama_reachable BOOL,
  generated_at TIMESTAMP)
\end{verbatim}

%% ════════════════════════════════════════════════════════════════════════════
\section{Environment Variables Reference}
\label{app:env}
%% ════════════════════════════════════════════════════════════════════════════

\begin{table}[H]
\centering
\caption{Complete Environment Variable Reference}
\label{tab:app_env}
\begin{tabular}{@{}p{3.5cm}p{1.8cm}p{2.2cm}@{}}
\toprule
\textbf{Variable} & \textbf{Default} & \textbf{Notes} \\
\midrule
\texttt{DATABASE\_URL}         & (required)    & asyncpg DSN \\
\texttt{REDIS\_URL}            & (required)    & redis:// DSN \\
\texttt{MQTT\_BROKER\_HOST}    & mosquitto     & Container name \\
\texttt{MQTT\_BROKER\_PORT}    & 1883          & TCP port \\
\texttt{ROOM\_ID}              & room1         & Must match ESP32 config \\
\texttt{MOCK\_MODE}            & false         & true: no ESP32 needed \\
\texttt{FACE\_RECOGNITION\_ENABLED} & false   & true: RPi + camera only \\
\texttt{SECRET\_KEY}           & changeme…     & MUST replace in prod \\
\texttt{REQUIRE\_AUTH}         & false         & true in production \\
\texttt{ACCESS\_TOKEN\_EXPIRE\_MINUTES} & 480 & JWT TTL \\
\texttt{TEMP\_AC\_ON\_THRESHOLD}  & 28         & °C \\
\texttt{TEMP\_AC\_OFF\_THRESHOLD} & 22         & °C \\
\texttt{AIR\_QUALITY\_ALERT\_THRESHOLD} & 500 & ppm \\
\texttt{FACE\_RECOGNITION\_THRESHOLD} & 0.40  & Cosine distance \\
\texttt{RECOGNITION\_FPS}      & 2             & Camera FPS \\
\texttt{AT\_RISK\_THRESHOLD}   & 0.70          & Fraction \\
\texttt{FORECAST\_WINDOW}      & 8             & Past sessions \\
\texttt{ATTENDANCE\_CYCLE\_DURATION} & 60     & Seconds/cycle \\
\texttt{OLLAMA\_BASE\_URL}     & http://ollama:11434 & Internal URL \\
\texttt{OLLAMA\_MODEL}         & phi3:mini     & Model tag \\
\texttt{MOODLE\_URL}           & (optional)    & LMS base URL \\
\texttt{MOODLE\_TOKEN}         & (optional)    & REST API token \\
\bottomrule
\end{tabular}
\end{table}

%% ════════════════════════════════════════════════════════════════════════════
\section{Hardware Wiring Diagram}
\label{app:wiring}
%% ════════════════════════════════════════════════════════════════════════════

The complete ESP32 GPIO assignment is shown in Table~\ref{tab:app_gpio}.
Figure~\ref{fig:esp32_wiring} (Section~\ref{subsec:iot_layer}) provides
the full wiring diagram.

\begin{table}[H]
\centering
\caption{ESP32 GPIO Pin Assignment}
\label{tab:app_gpio}
\begin{tabular}{@{}p{1.0cm}p{1.5cm}p{2.2cm}p{3.0cm}@{}}
\toprule
\textbf{GPIO} & \textbf{Dir.} & \textbf{Component} & \textbf{Signal} \\
\midrule
4  & In  & DHT21        & Temp/humidity (1-wire) \\
34 & In  & MQ-135       & Air quality ADC (0--4095) \\
35 & In  & ACP014       & Sound detect (digital) \\
21 & I2C & LCD 16×2     & I2C SDA \\
22 & I2C & LCD 16×2     & I2C SCL \\
26 & Out & Relay CH1/LLC & AC relay control \\
27 & Out & Relay CH2/LLC & Lighting relay control \\
14 & Out & Relay CH3/LLC & Spare \\
12 & Out & Relay CH4/LLC & Spare \\
\bottomrule
\end{tabular}
\end{table}

GPIO pins 34 and 35 are input-only and have no internal pull-up or
pull-down resistors.  The logic-level converter (3.3\,V $\leftrightarrow$
5\,V) is mandatory between the ESP32 outputs and the relay IN pins.

%% ════════════════════════════════════════════════════════════════════════════
\section{Docker Compose Configuration Summary}
\label{app:docker}
%% ════════════════════════════════════════════════════════════════════════════

\begin{table}[H]
\centering
\caption{Docker Compose Services}
\label{tab:app_docker}
\begin{tabular}{@{}p{1.8cm}p{2.2cm}p{0.8cm}p{2.7cm}@{}}
\toprule
\textbf{Service} & \textbf{Image} & \textbf{Port} & \textbf{Volume} \\
\midrule
postgres    & postgres:15-alpine    & 5432 & postgres\_data \\
redis       & redis:7-alpine        & 6379 & redis\_data \\
mosquitto   & eclipse-mosquitto:2   & 1883 & mosquitto\_data \\
backend     & custom (slim)         & 8000 & — \\
frontend    & custom (nginx)        & 3000 & — \\
ollama      & ollama/ollama:latest  & 11434 & ollama\_data \\
moodle*     & bitnami/moodle:4      & 8080 & moodle\_data \\
mariadb*    & mariadb:10.6          & 3306 & mariadb\_data \\
\bottomrule
\multicolumn{4}{l}{* Optional; activated with \texttt{--profile moodle}}
\end{tabular}
\end{table}

Quick-start commands:
\begin{verbatim}
# Full stack (with ESP32)
docker compose up -d

# Without ESP32 (mock sensors)
MOCK_MODE=true docker compose up -d

# With Moodle LMS
docker compose --profile moodle up -d

# Seed demo data
docker compose exec backend python seed.py
\end{verbatim}

%% ════════════════════════════════════════════════════════════════════════════
\section{Notion Project Board}
\label{app:notion}
%% ════════════════════════════════════════════════════════════════════════════

Figure~\ref{fig:notion_appendix} shows the Notion task board used for
project tracking across all 22 development phases.  Each task was
assigned to a team member, estimated in hours, and marked complete on
the day it was merged to the main branch.

\begin{figure}[H]
  \centering
  \includegraphics[width=\linewidth]{screenshot_notion}
  \caption{Notion project board: 22-phase delivery tracking with
    assignee and completion status.}
  \label{fig:notion_appendix}
\end{figure}
